\documentclass[12pt,a4paper]{article}

% --- Packages ---
\usepackage[utf8]{inputenc}
\usepackage[margin=2.5cm]{geometry}
\usepackage{amsmath,amssymb}
\usepackage{graphicx}
\usepackage{booktabs}
\usepackage{caption}
\usepackage{microtype}
\usepackage{cite}
\usepackage{hyperref}

\hypersetup{
    colorlinks=true,
    linkcolor=blue,
    citecolor=blue,
    urlcolor=blue
}

% --- Document Information ---
\title{\textbf{Judul Makalah / Laporan Ilmiah}}
\author{
    \textbf{Nama Penulis} \\
    NIM: 1234567890 \\
    Program Studi Teknik Informatika \\
    \texttt{email@example.com}
}
\date{\today}

\begin{document}

\maketitle

\begin{abstract}
Ini adalah bagian abstrak. Berikan ringkasan singkat mengenai latar belakang permasalahan, tujuan penelitian, metodologi yang diusulkan, hasil utama yang diperoleh, serta kesimpulan penting secara padat dan jelas.
\end{abstract}

\textbf{Kata Kunci:} LaTeX, Dev Container, Metodologi Penelitian, Reproducibility.

\section{Pendahuluan}
Dokumen ini disusun sebagai template penulisan ilmiah dengan dukungan lingkungan terisolasi menggunakan VS Code Dev Container dan TeX Live.

Salah satu tantangan dalam kolaborasi penulisan dokumen berbasis \LaTeX{} adalah perbedaan versi compiler dan paket di mesin masing-masing. Dengan Dev Container, seluruh dependensi sudah terinstalasi seragam di dalam Docker container \cite{knuth1984texbook}.

\section{Tinjauan Pustaka}
Bagian ini memuat teori dasar dan penelitian terdahulu yang relevan. Sebagai contoh, arsitektur transformator diperkenalkan oleh Vaswani dkk. \cite{vaswani2017attention}, yang berakar pada teori informasi Shannon \cite{shannon1948mathematical}.

\section{Metodologi Penelitian}
Uraikan langkah-langkah penelitian, arsitektur sistem, atau pendekatan matematis di bagian ini. Contoh rumus matematika:
\begin{equation}
    f(x) = \frac{1}{\sqrt{2\pi\sigma^2}} \exp\left(-\frac{(x - \mu)^2}{2\sigma^2}\right)
\end{equation}

\subsection{Tabel dan Data}
Contoh pembuatan tabel rapi dengan paket \texttt{booktabs}:
\begin{table}[htbp]
    \centering
    \caption{Contoh Hasil Evaluasi Model}
    \label{tab:evaluasi}
    \begin{tabular}{lccc}
        \toprule
        \textbf{Model} & \textbf{Precision} & \textbf{Recall} & \textbf{F1-Score} \\
        \midrule
        Baseline       & 0.81               & 0.79            & 0.80              \\
        Proposed Model & \textbf{0.92}      & \textbf{0.89}   & \textbf{0.90}     \\
        \bottomrule
    \end{tabular}
\end{table}

\section{Hasil dan Pembahasan}
Jelaskan interpretasi dari data dan hasil eksperimen yang didapatkan pada Tabel \ref{tab:evaluasi}.

\section{Kesimpulan}
Tuliskan kesimpulan penelitian serta saran untuk pengembangan di masa mendatang.

% --- Referensi ---
\bibliographystyle{IEEEtran}
\bibliography{references}

\end{document}
